\documentclass[a4paper,11pt]{ltjsarticle}
\usepackage[haranoaji]{luatexja-preset}
\usepackage{amsmath,amssymb}
\usepackage[top=12mm,bottom=10mm,left=14mm,right=14mm]{geometry}
\usepackage{array}
\usepackage{tikz}
\usetikzlibrary{calc}
\pagestyle{empty}
\setlength{\parindent}{0pt}
\newlength{\qcol}\setlength{\qcol}{0.47\textwidth}
\newlength{\qgap}
\newlength{\anscolw}\setlength{\anscolw}{0.30\textwidth}
\newlength{\ansh}
\newcommand{\Q}[2]{\parbox[t]{\qcol}{\raggedright (#1)\hspace{0.6em}$\displaystyle #2$}}
\newcommand{\QR}[4]{\Q{#1}{#2}\hfill\Q{#3}{#4}\par\vspace{\qgap}}
\newcommand{\anscell}[1]{\rule[-\ansdrop]{0pt}{\ansh}(#1)}
\newlength{\ansdrop}

\begin{document}
{\large\bfseries 計算問題　27}\hfill 名前(\underline{\hspace{32mm}})\par\vspace{3mm}
■（1）〜（16）の計算をしなさい。（17）〜（21）は方程式を解きなさい。\par\vspace{3mm}
\setlength{\qgap}{5.4mm}
\QR{1}{7\div3\times(-6)}{2}{(-21)\div(+3)}
\QR{3}{\dfrac{x-7}{3}\times6}{4}{(21a-6)\div3}
\QR{5}{8\times\{3-(-2+1)\times2\}}{6}{(-14)+(-6)-(-7)}
\QR{7}{-7^2+(-3)^2-2}{8}{(7a+3)+(2a-8)}
\QR{9}{7(3x-2)-(2x-8)}{10}{15-(-3)\div\left(-\dfrac{3}{2}\right)}
\QR{11}{(3x-7)-(8x+2)}{12}{\dfrac{7}{3}(3x-24)-\dfrac{2}{8}(8x+56)}
\QR{13}{\dfrac{1}{7}(21x-14)+\dfrac{1}{8}(16x+24)}{14}{-0.7-(-2.3)-1.8}
\QR{15}{(-11)\times(-5)}{16}{\dfrac{7}{3}\div2+\dfrac{8}{3}}
\QR{17}{7x-3=4}{18}{7(x+3)=8x+18}
\QR{19}{0.7x+0.3=1.7}{20}{\dfrac{x+5}{3}=\dfrac{x-1}{2}}
\vspace{1mm}
\Q{21}{7(x-2)-3(x+1)=10x-13}\par\vspace{3mm}
\setlength{\ansh}{9.5mm}\setlength{\ansdrop}{6mm}%
\begin{center}\begin{tabular}{|p{\anscolw}|p{\anscolw}|p{\anscolw}|}\hline
\anscell{1} & \anscell{2} & \anscell{3} \\\hline
\anscell{4} & \anscell{5} & \anscell{6} \\\hline
\anscell{7} & \anscell{8} & \anscell{9} \\\hline
\anscell{10} & \anscell{11} & \anscell{12} \\\hline
\anscell{13} & \anscell{14} & \anscell{15} \\\hline
\anscell{16} & \anscell{17} & \anscell{18} \\\hline
\anscell{19} & \anscell{20} & \anscell{21} \\\hline
\end{tabular}\end{center}

\end{document}
